\section{A trace certificate and quantitative finite-power stability}
This section gives a second proof of tensor destabilization, using the
trace of the entire tangent Hessian. It produces a two-column rotation
in the lifted dimension and a bound on the energy of any matrix whose
specified tensor power is locally minimizing. Neither conclusion
supplies a descent in the original dimension.

Let $U\in\mathbb R^{d\times n}$ have orthonormal rows, $1\le d\le n$,
$n\ge2$, and let $p>2$. Write
\[
 q=\sum_{ij}|u_{ij}|^p,\qquad r=\sum_{ij}|u_{ij}|^{p-2},
 \qquad \rho=\frac{r}{nq}.
\]
By strict Chebyshev inequality, $0<\rho\le1$, with equality exactly
when every entry has magnitude $n^{-1/2}$.

\begin{theorem}[Finite-power trace test]\label{thm:trace-test}
If $U^{\otimes k}$ is a local minimum of the entrywise $p$-energy on
the row-Stiefel manifold, then
\begin{equation}\label{eq:trace-test}
 (p-1)\rho^k\ \ge\ 1+\frac{p-2}{n^k}.
\end{equation}
In particular, a nonflat $U$ has a strictly negative second variation
along a single two-column rotation of $U^{\otimes k}$ whenever
$k>\log(p-1)/\log(1/\rho)$. The second-variation conclusion does not
require that $U$ be critical.
\end{theorem}

\begin{proof}
For any real row-orthonormal $V$ with $N$ columns, let
$Q=\sum |v_{ij}|^p$ and $R=\sum |v_{ij}|^{p-2}$. For
$X_{ab}=E_{ab}-E_{ba}$, direct differentiation of
$F_p(Ve^{tX_{ab}})$ gives
\begin{align*}
 \frac1p\frac{d^2}{dt^2}F_p(Ve^{tX_{ab}})\bigg|_{t=0}
 ={}&(p-1)\sum_i\bigl(|v_{ia}|^{p-2}v_{ib}^2
               +|v_{ib}|^{p-2}v_{ia}^2\bigr)\\
 &-\sum_i\bigl(|v_{ia}|^p+|v_{ib}|^p\bigr).
\end{align*}
Summing over $a<b$ and using $\sum_bv_{ib}^2=1$ yields
\begin{equation}\label{eq:curvature-trace}
 \frac1p\sum_{a<b}\frac{d^2}{dt^2}F_p(Ve^{tX_{ab}})\bigg|_{t=0}
   =(p-1)R-(N+p-2)Q.
\end{equation}
For $V=U^{\otimes k}$ we have $N=n^k$, $Q=q^k$, $R=r^k$.
A local minimum has nonnegative second derivative on every such
curve, which gives \eqref{eq:trace-test}. If
$(p-1)\rho^k<1$, the displayed sum is strictly negative; at least one
of its summands is strictly negative. Smoothness at zero is valid
because $p>2$.
\end{proof}

The column pair can be checked without forming the tensor matrix.
Define $c_a=\sum_i|u_{ia}|^p$ and
$D_{ab}=\sum_i|u_{ia}|^{p-2}u_{ib}^2$. For distinct column words
$\mathbf a,\mathbf b\in\{1,\ldots,n\}^k$, the second derivative divided
by $p$ is exactly
\[
 (p-1)\left(\prod_{j=1}^kD_{a_jb_j}
             +\prod_{j=1}^kD_{b_ja_j}\right)
       -\prod_{j=1}^kc_{a_j}-\prod_{j=1}^kc_{b_j}.
\]
This is an explicit certificate formula, not a claim of a polynomial
time search algorithm. A two-column rotation of a tensor matrix
usually does not factor as rotations of its original factors.

\begin{corollary}[Dimension-independent quartic stability bound]
\label{cor:quartic-rate}
Let $e=nF_4(U)/d\ge1$. If $U^{\otimes k}$ is a local minimum, then
\[
 e\le \left(\frac{3}{1+2/n^k}\right)^{1/k}<3^{1/k}.
\]
Furthermore
\[
 \sum_{ij}\left(u_{ij}^2-\frac1n\right)^2
      =\frac dn(e-1),\qquad
 \frac1d\min_{H\in\{\pm1\}^{d\times n}}
       \left\|U-\frac H{\sqrt n}\right\|_F^2
      \le e-1<3^{1/k}-1.
\]
Thus the normalized squared distance to a flat sign array is
$O(1/k)$, uniformly in the matrix dimensions.
\end{corollary}

\begin{proof}
At $p=4$ we have $r=d$, hence $\rho=1/e$; apply
\eqref{eq:trace-test}. The first distance identity follows by expanding
the square and using $\sum u_{ij}^2=d$. Choose each sign of $H$ to
match $u_{ij}$ (arbitrarily at zero) and use
\[
 (|u_{ij}|-n^{-1/2})^2
 =\frac{(u_{ij}^2-n^{-1})^2}{(|u_{ij}|+n^{-1/2})^2}
 \le n(u_{ij}^2-n^{-1})^2.
\]
Finally $3^{1/k}-1\sim(\log3)/k$.
\end{proof}

The closest flat sign array in this estimate need not have orthogonal
rows. The estimate therefore proves neither a nearby Hadamard matrix
nor Hadamard existence. It quantifies a necessary condition on local
minima in tensor powers.

\subsection{Complex flatness and collision entropy}
For complex row-orthonormal $U$, put again $e=nF_4(U)/d$. The complex
trace identity gives the slightly stronger necessary condition
\[
 e\le\left(\frac{2}{1+1/n^k}\right)^{1/k}<2^{1/k}
\]
when $U^{\otimes k}$ is locally minimizing. The same variance identity
holds with $u_{ij}^2$ replaced by $|u_{ij}|^2$; replacing signs by
arbitrary unit complex phases gives the corresponding distance bound.
This follows by choosing the phase of each entry to match $u_{ij}$.

There is a direct information-theoretic interpretation of both bounds.
Define a probability distribution on the $dn$ entry positions by
$P_{ij}=|u_{ij}|^2/d$, and let $Q_{ij}=1/(dn)$ be uniform. Its
order-two R\'enyi divergence from uniform, here defined explicitly, is
\[
 D_2(P\Vert Q)=\log\sum_{ij}\frac{P_{ij}^2}{Q_{ij}}
             =\log e.
\]
The entry distribution of $U^{\otimes k}$ is $P^{\otimes k}$, so
$D_2(P^{\otimes k}\Vert Q^{\otimes k})=kD_2(P\Vert Q)$.
Consequently finite-power local minimality forces
\[
 D_2(P\Vert Q)<\frac{\log3}{k}\quad\text{over }\mathbb R,
 \qquad
 D_2(P\Vert Q)<\frac{\log2}{k}\quad\text{over }\mathbb C.
\]
This expresses the curvature test as an entropy-deficit bound. It
does not assert that an arbitrary nearly uniform entry distribution
can be realized by a flat orthogonal matrix of the original size.
